% !TeX spellcheck = en_US
\documentclass[french]{yLectureNote}

\title{Outils Mathématiques 2}
\subtitle{Physique}
\author{Paulhenry Saux}
\date{\today}
\yLanguage{Français}

\professor{F.Pettinari}
\usepackage{graphicx}%----pour mettre des images
\usepackage[utf8]{inputenc}%---encodage
\usepackage{geometry}%---pour modifier les tailles et mettre a4paper
%\usepackage{awesomebox}%---pour les boites d'exercices, de pbq et de croquis ---d\'esactiv\'e pour les TP de PC
\usepackage{tikz}%---pour deiffner + d\'ependance de chemfig
% \usepackage{tabularx}%---pour dimensionner automatiquement les tableaux avec variable X
\usepackage{awesomebox}%---Pour les boites info, danger et autres
\usepackage{menukeys}%---Pour deiffner les touches de Calculatrice
\usepackage{fancyhdr}%---pour les en-t\^ete personnalis\'ees
\usepackage{blindtext}%---pour les liens
\usepackage{hyperref}%---pour les liens (\`a mettre en dernier)
\usepackage{caption}%---pour la francisation de la l\'egende table vers Tableau
\usepackage{pifont}
\usepackage{array}%---pour les tableaux
\usepackage{lipsum}
\usepackage{yFlatTable}
\usepackage{multicol}
\newcommand{\Lim}[1]{\lim\limits_{\substack{#1}}\:}
\renewcommand{\vec}{\overrightarrow}
\newcommand{\N}[0]{\mathbb{N}}
\newcommand{\dd}{\mathrm{d}}
\newcommand{\norm}[1]{||\vec{#1}||}
\newcommand{\fo}{\psi(\vec{r},t)}
\newcommand{\foe}{\psi(\vec{r},t)\*}
\newcommand{\HH}{\hat{H}}
\newcommand{\hb}{\hbar}
\newcommand{\lap}{\nabla^2}
\newcommand{\lapcc}{\frac{\partial^2 }{\partial x^2}+\frac{\partial^2 }{\partial y^2}+\frac{\partial^2 }{\partial z^2}}
\newcommand{\mpsi}{\(\psi\)}
\begin{document}
\setcounter{chapter}{2}
\chapter{Calcul vectoriel}
\section{Paramétrisation}
\subsection{Segment du point $A=(1,0)$ au point $B=(0,1)$}
La paramétrisation d'un segment de droite joignant $A$ à $B$ est donnée par la formule :
$$\vec{r}(t) = \vec{OA} + t \cdot \vec{AB}, \quad t \in [0, 1]$$

1.  Vecteur directeur $\vec{AB}$ :
    $$\vec{AB} = \vec{OB} - \vec{OA} = \begin{pmatrix} 0 - 1 \\ 1 - 0 \end{pmatrix} = \begin{pmatrix} -1 \\ 1 \end{pmatrix}$$

2.  Paramétrisation $\vec{r}(t)$ :
    $$\vec{r}(t) = \begin{pmatrix} 1 \\ 0 \end{pmatrix} + t \begin{pmatrix} -1 \\ 1 \end{pmatrix} = \begin{pmatrix} 1 - t \\ t \end{pmatrix}$$
    $$\vec{r}(t) = (1 - t) \vec{e}_x + t \vec{e}_y, \quad t \in [0, 1]$$

3.  Dérivée du vecteur position $\frac{d\vec{r}}{dt}$ :
    $$\frac{d\vec{r}}{dt} = \begin{pmatrix} \frac{d}{dt}(1 - t) \\ \frac{d}{dt}(t) \end{pmatrix}$$
    $$\frac{d\vec{r}}{dt} = - \vec{e}_x + \vec{e}_y = \begin{pmatrix} -1 \\ 1 \end{pmatrix}$$

\subsection{Segment du point $A=(0,3)$ au point $B=(4,1)$ }

On utilise la même formule $\vec{r}(t) = \vec{A} + t \cdot \vec{AB}, \quad t \in [0, 1]$.

1.  Vecteur directeur $\vec{AB}$ :
    $$\vec{AB} = \vec{B} - \vec{A} = \begin{pmatrix} 4 - 0 \\ 1 - 3 \end{pmatrix} = \begin{pmatrix} 4 \\ -2 \end{pmatrix}$$

2.  Paramétrisation $\vec{r}(t)$ :
    $$\vec{r}(t) = \begin{pmatrix} 0 \\ 3 \end{pmatrix} + t \begin{pmatrix} 4 \\ -2 \end{pmatrix} = \begin{pmatrix} 4t \\ 3 - 2t \end{pmatrix}$$
    $$\vec{r}(t) = 4t \vec{e}_x + (3 - 2t) \vec{e}_y, \quad t \in [0, 1]$$

3.  Dérivée du vecteur position $\frac{d\vec{r}}{dt}$ :
    $$\frac{d\vec{r}}{dt} = \begin{pmatrix} \frac{d}{dt}(4t) \\ \frac{d}{dt}(3 - 2t) \end{pmatrix}$$
    $$\frac{d\vec{r}}{dt} = 4 \vec{e}_x - 2 \vec{e}_y = \begin{pmatrix} 4 \\ -2 \end{pmatrix}$$
\subsection{Arc de cercle de centre $C=(-1,0)$, d'extrémités $A=(0,0)$ et $B=(-1,1)$ de A vers B.}

1.  Détermination du Rayon ($R$) et du Centre ($C$) : Centre : $C(-1, 0)$ et le rayon $R = CA = \sqrt{(0 - (-1))^2 + (0 - 0)^2} = 1$.

2.  Paramétrisation :
    Pour un cercle de centre $(x_C, y_C)$ et rayon $R$, la paramétrisation est :
    $$\begin{cases} x(t) = x_C + R \cos(t) \\ y(t) = y_C + R \sin(t) \end{cases} \implies \begin{cases} x(t) = -1 + \cos(t) \\ y(t) = \sin(t) \end{cases}$$

3.  Détermination de l'intervalle $t$ :
    Le paramètre $t$ représente l'angle $\theta$ en coordonnées polaires par le centre $C$.
    Point A ($0, 0$) :
        $$\begin{cases} 0 = -1 + \cos(t) \implies \cos(t) = 1 \\ 0 = \sin(t) \implies \sin(t) = 0 \end{cases} \implies t_{initial} = 0$$
    Point B ($-1, 1$) :
        $$\begin{cases} -1 = -1 + \cos(t) \implies \cos(t) = 0 \\ 1 = \sin(t) \implies \sin(t) = 1 \end{cases} \implies t_{final} = \frac{\pi}{2}$$

Remarque : on pouvait évidement trouver t à l'oeil, mais la justification rigoureuse est celle ci-dessus.

    L'arc de cercle dirigé de $A$ vers $B$ est donc :
    $$\vec{r}(t) = (-1 + \cos(t)) \vec{e}_x + \sin(t) \vec{e}_y, \quad t \in [0, \frac{\pi}{2}]$$

4.  Dérivée du vecteur position $\frac{\dd\vec{r}}{dt}$ :
    $$\frac{\dd\vec{r}}{dt} = \begin{pmatrix} \frac{d}{dt}(-1 + \cos(t)) \\ \frac{d}{dt}(\sin(t)) \end{pmatrix}$$
    $$\frac{\dd\vec{r}}{dt} = - \sin(t) \vec{e}_x + \cos(t) \vec{e}_y = \begin{pmatrix} -\sin(t) \\ \cos(t) \end{pmatrix}$$
\section{Intégrales de champ vectoriel}
\subsection{Valeur moyenne du champ de vitesse}
Le champ de vitesse dans un tube de rayon $R$ est donné par $\vec{v}(\vec{r})=\vec{v}_{max}(1-\rho^{2}/R^{2})\vec{e}_z$ en coordonnées cylindriques.

La valeur moyenne $\langle \vec{v} \rangle$ sur une section $\mathcal{S}$ (un disque de rayon $R$) est donnée par :
$$\langle \vec{v} \rangle = \frac{1}{\text{Aire}(\mathcal{S})} \iint_{\mathcal{S}} \vec{v} \, dS$$

1.  Aire de la section : $\text{Aire}(\mathcal{S}) = \pi R^2$.
2.  Calcul de l'intégrale : On utilise les coordonnées polaires sur le disque ($dS = \rho \, d\rho \, d\varphi$) :
\begin{flalign*}
\iint_{\mathcal{S}} \vec{v} \, dS &= \int_{0}^{2\pi} d\varphi \cdot \int_{0}^{R} \vec{v}_{max} \left(1 - \frac{\rho^{2}}{R^{2}}\right) \rho \, d\rho \, \vec{e}_z\\
    &= 2\pi \vec{v}_{max} \vec{e}_z \cdot \int_{0}^{R} \left(\rho - \frac{\rho^{3}}{R^{2}}\right) d\rho\\
    &= 2\pi \vec{v}_{max} \vec{e}_z \cdot \left[ \frac{\rho^{2}}{2} - \frac{\rho^{4}}{4R^{2}} \right]_{0}^{R}\\
    &= 2\pi \vec{v}_{max} \vec{e}_z \cdot \left( \frac{R^{2}}{2} - \frac{R^{2}}{4} \right) \\
    &= \frac{\pi R^{2}}{2} \vec{v}_{max} \vec{e}_z
\end{flalign*}



3.  Valeur moyenne :
    $$\langle \vec{v} \rangle = \frac{1}{\pi R^2} \cdot \left( \frac{\pi R^{2}}{2} \vec{v}_{max} \vec{e}_z \right) = \frac{1}{2} \vec{v}_{max} \vec{e}_z$$
\subsection{Valeur moyenne du vecteur $\vec{OM}$ dans le demi-disque $\mathcal{D}$}

Le domaine est $\mathcal{D}=\{(x,y),x^{2}+y^{2}<R^{2},y>0\}$. Le vecteur position est $\vec{OM} = x \vec{e}_x + y \vec{e}_y$.

$$\langle \vec{OM} \rangle = \frac{1}{\text{Aire}(\mathcal{D})} \iint_{\mathcal{D}} \vec{OM} \, dS$$

1.  Aire de $\mathcal{D}$ : $\text{Aire}(\mathcal{D}) = \frac{1}{2} \pi R^2$.
2.  Calcul de l'intégrale : On utilise les coordonnées polaires ($\rho \in [0, R]$, $\varphi \in [0, \pi]$).
    $$\iint_{\mathcal{D}} \vec{OM} \, dS = \int_{0}^{\pi} \int_{0}^{R} (\rho \cos\varphi \vec{e}_x + \rho \sin\varphi \vec{e}_y) \rho \, d\rho \, d\varphi$$
    La composante en $\vec{e}_x$ est nulle car $\int_{0}^{\pi} \cos\varphi \, d\varphi = 0$.

    La composante en $\vec{e}_y$ :
    \begin{flalign*}
    \int_{0}^{\pi} \sin\varphi \, d\varphi \cdot \int_{0}^{R} \rho^{2} \, d\rho &= [-\cos\varphi]_{0}^{\pi} \cdot \left[ \frac{\rho^{3}}{3} \right]_{0}^{R}\\
    &= (1 - (-1)) \cdot \frac{R^{3}}{3}\\
    &= 2 \cdot \frac{R^{3}}{3} \\
    &= \frac{2R^{3}}{3}
    \end{flalign*}
    Donc, $\iint_{\mathcal{D}} \vec{OM} \, dS = \frac{2R^{3}}{3} \vec{e}_y$.
3.  Valeur moyenne :
    $$\langle \vec{OM} \rangle = \frac{1}{\frac{1}{2} \pi R^2} \cdot \frac{2R^{3}}{3} \vec{e}_y = \frac{4R^{3}}{3\pi R^2} \vec{e}_y = \langle \vec{OM} \rangle = \frac{4R}{3\pi} \vec{e}_y$$

\subsection{Circulation du poids $\vec{P}=-mg\vec{e}_{z}$}

Le poids $\vec{P}$ est un champ conservatif (dérivant du potentiel $V = -mg z$). La circulation le long d'un chemin $\mathcal{C}$ de $A$ à $B$ est égale à l'opposé de la variation d'énergie potentielle :
$$W_{A \to B} = \int_{\mathcal{C}} \vec{P} \cdot d\vec{r}$$

1.  Circulation sur le segment de $A$ à $B$ :
    $$W_{A \to B} = \int_{\mathcal{C}} (-mg\vec{e}_{z}) \cdot d\vec{r}$$
    Comme $\vec{e}_z \cdot d\vec{r} = dz$, on a :
    $$W_{A \to B} = \int_{z_A}^{z_B} -mg \, dz = -mg [z]_{z_A}^{z_B}$$
    $$\mathbf{W_{A \to B} = -mg(z_{B}-z_{A})}$$

2.  Circulation sur le chemin $A \to C \to B$ :
    Puisque le champ est conservatif, la circulation est indépendante du chemin et ne dépend que des extrémités.
    $$W_{A \to C \to B} = W_{A \to C} + W_{C \to B} = -mg(z_C - z_A) + (-mg(z_B - z_C))$$
    $$W_{A \to C \to B} = -mg(z_C - z_A + z_B - z_C) = \mathbf{-mg(z_{B}-z_{A})}$$

\subsection{Circulation de la force $\vec{F}(x,y)=\frac{F_{0}}{R}(-y\vec{e}_{x}+x\vec{e}_{y})$ sur le cercle $\mathcal{C}$}

Le cercle $\mathcal{C}$ a pour rayon $R$ et centre $O$, orienté positivement.

    Paramétrisation : $\vec{r}(t) = R \cos t \vec{e}_x + R \sin t \vec{e}_y$, $t \in [0, 2\pi]$.

    Déplacement : $d\vec{r} = (-R \sin t \vec{e}_x + R \cos t \vec{e}_y) dt$.

    Force sur $\mathcal{C}$ : $\vec{F}(\vec{r}(t)) = \frac{F_{0}}{R} (-R \sin t \vec{e}_{x} + R \cos t \vec{e}_{y}) = F_{0} (-\sin t \vec{e}_{x} + \cos t \vec{e}_{y})$.

    Produit scalaire :
        $$\vec{F} \cdot d\vec{r} = F_{0} [ (-\sin t)(-R \sin t) + (\cos t)(R \cos t) ] dt = F_{0} R (\sin^{2} t + \cos^{2} t) dt = F_{0} R \, dt$$

    Circulation :
        $$\oint_{\mathcal{C}} \vec{F} \cdot d\vec{r} = \int_{0}^{2\pi} F_{0} R \, dt = 2\pi R F_{0}$$


\subsection{Flux}
On cherche à calculer le flux de $\vec{F}(x,y,z)=\frac{1}{1+(z/l)^{2}}(x\vec{e}_{x}+y\vec{e}_{y}+z\vec{e}_{z})$ à travers la surface latérale $\mathcal{S}$ du cylindre infini d'axe $(O, \vec{e}_z)$ et de rayon $R$.

1.  Élément de surface orienté ($d\vec{S}$) : Sur le cylindre de rayon $R$, $d\vec{S} = \vec{e}_\rho \, dS = \vec{e}_\rho \, R \, d\varphi \, dz$.

2.  Produit scalaire $\vec{F} \cdot d\vec{S}$ :
    $$\vec{F} \cdot d\vec{S} = \frac{1}{1+(z/l)^{2}} (x\vec{e}_{x}+y\vec{e}_{y}+z\vec{e}_{z}) \cdot \vec{e}_\rho \, R \, d\varphi \, dz$$

    En coordonnées cylindriques, $x\vec{e}_{x}+y\vec{e}_{y} = \rho \vec{e}_\rho$. Sur le cylindre, $\rho=R$.
    $$(x\vec{e}_{x}+y\vec{e}_{y}+z\vec{e}_{z}) = R\vec{e}_{\rho} + z\vec{e}_{z}$$
    $$(R\vec{e}_{\rho} + z\vec{e}_{z}) \cdot \vec{e}_\rho = R(\vec{e}_{\rho}\cdot\vec{e}_{\rho}) + z(\vec{e}_{z}\cdot\vec{e}_{\rho}) = R$$
    $$\vec{F} \cdot d\vec{S} = \frac{R}{1+(z/l)^{2}} \, R \, d\varphi \, dz = \frac{R^2}{1+(z/l)^{2}} \, d\varphi \, dz$$
3.  Calcul du Flux : L'intégration se fait sur $z \in ]-\infty, +\infty[$ et $\varphi \in [0, 2\pi]$.
    $$\Phi = R^2 \int_{0}^{2\pi} d\varphi \cdot \int_{-\infty}^{+\infty} \frac{1}{1+(z/l)^{2}} \, dz$$
    $$\Phi = 2\pi R^2 \cdot (l [\arctan(z/l)]_{-\infty}^{+\infty}) = 2\pi R^2 \cdot (l (\frac{\pi}{2} - (-\frac{\pi}{2}))) =  2\pi^2 l R^2$$
\subsection{Flux de $\vec{F}$ champ de vecteur constant}

Soit $\vec{F}$ un champ constant, $S$ l'aire de la surface plane, et $\vec{n}$ son vecteur normal unitaire constant.

Le flux est donné par :
$$\Phi = \iint_{\mathcal{S}} \vec{F} \cdot d\vec{S}$$
Comme $\vec{F}$ est constant et $d\vec{S} = \vec{n} dS$ avec $\vec{n}$ constant :
$$\Phi = (\vec{F} \cdot \vec{n}) \iint_{\mathcal{S}} dS$$
$$\mathbf{\Phi = (\vec{F} \cdot \vec{n}) S}$$
\section{Calcul de gradient}
\subsection{1. Gradient du potentiel quadratique et force associée}

L'énoncé donne le potentiel scalaire $V(x,y)=k(x^{2}+y^{2})/2$ avec $k>0$.

Le gradient de $V$ en coordonnées cartésiennes $(x,y)$ est donné par:
$$
\vec{\nabla}V = \frac{\partial V}{\partial x}\vec{e}_{x} + \frac{\partial V}{\partial y}\vec{e}_{y}
$$
Calculons les dérivées partielles:
$\frac{\partial V}{\partial x} = \frac{\partial}{\partial x} \left( \frac{k}{2}(x^{2}+y^{2}) \right) = \frac{k}{2}(2x) = kx$

$\frac{\partial V}{\partial y} = \frac{\partial}{\partial y} \left( \frac{k}{2}(x^{2}+y^{2}) \right) = \frac{k}{2}(2y) = ky$

Donc, le gradient de $V$ est:
$$
\vec{\nabla}V(x,y) = kx\vec{e}_{x} + ky\vec{e}_{y} = k(x\vec{e}_{x} + y\vec{e}_{y}) = k\vec{r}
$$
où $\vec{r} = x\vec{e}_{x} + y\vec{e}_{y}$ est le vecteur position.


La force $\vec{F}$ ressentie dans un potentiel $V$ est donnée par $\vec{F} = -\vec{\nabla}V$.
$$
\vec{F}(x,y) = -k(x\vec{e}_{x} + y\vec{e}_{y}) = -k\vec{r}
$$

Les courbes de niveau de $V(x,y)$ sont les courbes d'équation $V(x,y) = C$ (avec $C \ge 0$ une constante).
$$
\frac{k}{2}(x^{2}+y^{2}) = C \quad \Rightarrow \quad x^{2}+y^{2} = \frac{2C}{k} = R^2
$$
Les courbes de niveau sont des cercles centrés à l'origine $O$.

Le gradient $\vec{\nabla}V$ (et donc la force $\vec{F}$) est toujours orthogonal aux courbes de niveau et pointe vers les valeurs croissantes du potentiel (vers l'extérieur). La force $\vec{F}$ est dirigée vers le centre $O$.
\subsection{Gradient du potentiel gravitationnel}

On a la fonction scalaire $f(x,y,z)=1/\sqrt{x^{2}+y^{2}+z^{2}}$. En utilisant la notation $r = \sqrt{x^{2}+y^{2}+z^{2}}$, nous avons $f(\vec{r}) = 1/r$.


Pour calculer le gradient $\vec{\nabla}f$, il est souvent plus simple d'utiliser la dérivée par rapport à $r$ et l'identité $\vec{\nabla}g(r) = g'(r)\vec{e}_{r}$, où $\vec{e}_{r} = \vec{r}/r$ est le vecteur unitaire radial.

Ici, $f(r) = 1/r$.
La dérivée de $f$ par rapport à $r$ est:
$$
f'(r) = \frac{d}{dr} \left(\frac{1}{r}\right) = -\frac{1}{r^{2}}
$$
Donc, le gradient est:
$$
\vec{\nabla}f = f'(r)\vec{e}_{r} = -\frac{1}{r^{2}}\vec{e}_{r}
$$
En coordonnées cartésiennes, rappelant que $\vec{e}_{r} = \frac{x\vec{e}_{x} + y\vec{e}_{y} + z\vec{e}_{z}}{r}$:
$$
\vec{\nabla}f(x,y,z) = -\frac{1}{r^{2}}\frac{x\vec{e}_{x} + y\vec{e}_{y} + z\vec{e}_{z}}{r} = -\frac{x\vec{e}_{x} + y\vec{e}_{y} + z\vec{e}_{z}}{(x^{2}+y^{2}+z^{2})^{3/2}}
$$

Le potentiel gravitationnel ressenti par une particule de masse $m$ au point $\vec{r}$ en présence d'une masse $M$ à l'origine est $V(\vec{r})=-\mathcal{G}mM/r$.
La force gravitationnelle $\vec{F}_{G}$ est donnée par:
$$
\vec{F}_{G} = -\vec{\nabla}V
$$
Puisque $V(\vec{r}) = -\mathcal{G}mM \cdot \left(\frac{1}{r}\right) = -\mathcal{G}mM \cdot f(\vec{r})$, nous avons:
$$
\vec{F}_{G} = -\vec{\nabla}\left(-\frac{\mathcal{G}mM}{r}\right) = \mathcal{G}mM \cdot \vec{\nabla}\left(\frac{1}{r}\right)
$$
En utilisant le résultat du gradient de $1/r$:
$$
\vec{F}_{G} = \mathcal{G}mM \cdot \left(-\frac{1}{r^{2}}\vec{e}_{r}\right) = -\frac{\mathcal{G}mM}{r^{2}}\vec{e}_{r}
$$
C'est la loi de la gravitation universelle de Newton. La force est attractive (dirigée par $-\vec{e}_{r}$ vers la masse $M$ à l'origine) et sa magnitude est $\frac{\mathcal{G}mM}{r^{2}}$.
\section{Changement de coordonnées}
\subsection{Coordonnées Polaires}
\subsubsection{Vérification}
Nous voulons exprimer $\vec{\nabla}f = \frac{\partial f}{\partial x}\vec{e}_{x} + \frac{\partial f}{\partial y}\vec{e}_{y}$ dans la base polaire $(\vec{e}_{\rho}, \vec{e}_{\varphi})$.

1. Relations

On sait que :
$$
x = \rho \cos \varphi \quad \text{et} \quad y = \rho \sin \varphi
$$

De plus, les vecteurs unitaires polaires $\vec{e}_{\rho}$ et $\vec{e}_{\varphi}$ s'expriment dans la base cartésienne $(\vec{e}_{x}, \vec{e}_{y})$ :
$$
\vec{e}_{\rho} = \cos \varphi \vec{e}_{x} + \sin \varphi \vec{e}_{y}
$$
$$
\vec{e}_{\varphi} = -\sin \varphi \vec{e}_{x} + \cos \varphi \vec{e}_{y}
$$

Pour trouver les composantes cartésiennes $\vec{e}_{x}$ et $\vec{e}_{y}$ en fonction des vecteurs polaires, nous inversons ces relations :
$$
\vec{e}_{x} = \cos \varphi \vec{e}_{\rho} - \sin \varphi \vec{e}_{\varphi}
$$
$$
\vec{e}_{y} = \sin \varphi \vec{e}_{\rho} + \cos \varphi \vec{e}_{\varphi}
$$

2. Relations entre les Dérivées Partielles

En utilisant la règle de la chaîne, nous exprimons les dérivées cartésiennes en fonction des dérivées polaires. Nous aurons besoin des dérivées suivantes (à partir des relations en A):

 $\frac{\partial x}{\partial \rho} = \cos \varphi$

 $\frac{\partial y}{\partial \rho} = \sin \varphi$

 $\frac{\partial x}{\partial \varphi} = -\rho \sin \varphi$

 $\frac{\partial y}{\partial \varphi} = \rho \cos \varphi$

Les relations de la règle de la chaîne sont :
$$
\frac{\partial f}{\partial \rho} = \frac{\partial f}{\partial x}\frac{\partial x}{\partial \rho} + \frac{\partial f}{\partial y}\frac{\partial y}{\partial \rho} = \frac{\partial f}{\partial x}\cos \varphi + \frac{\partial f}{\partial y}\sin \varphi \quad \mathbf{(1)}
$$
$$
\frac{\partial f}{\partial \varphi} = \frac{\partial f}{\partial x}\frac{\partial x}{\partial \varphi} + \frac{\partial f}{\partial y}\frac{\partial y}{\partial \varphi} = \frac{\partial f}{\partial x}(-\rho \sin \varphi) + \frac{\partial f}{\partial y}(\rho \cos \varphi) \quad \mathbf{(2)}
$$

Pour utiliser ces relations dans l'expression du gradient, nous divisons $\mathbf{(2)}$ par $\rho$ :
$$
\frac{1}{\rho}\frac{\partial f}{\partial \varphi} = -\frac{\partial f}{\partial x}\sin \varphi + \frac{\partial f}{\partial y}\cos \varphi \quad \mathbf{(3)}
$$

3. Dérivation du Gradient

Le gradient en coordonnées cartésiennes est :
$$
\vec{\nabla}f = \frac{\partial f}{\partial x}\vec{e}_{x} + \frac{\partial f}{\partial y}\vec{e}_{y}
$$

Remplaçons $\vec{e}_{x}$ et $\vec{e}_{y}$ par leurs expressions en base polaire :
$$
\vec{\nabla}f = \frac{\partial f}{\partial x}(\cos \varphi \vec{e}_{\rho} - \sin \varphi \vec{e}_{\varphi}) + \frac{\partial f}{\partial y}(\sin \varphi \vec{e}_{\rho} + \cos \varphi \vec{e}_{\varphi})
$$

Regroupons les termes selon les vecteurs de base polaire $\vec{e}_{\rho}$ et $\vec{e}_{\varphi}$ :
$$
\vec{\nabla}f = \left(\frac{\partial f}{\partial x}\cos \varphi + \frac{\partial f}{\partial y}\sin \varphi\right)\vec{e}_{\rho} + \left(-\frac{\partial f}{\partial x}\sin \varphi + \frac{\partial f}{\partial y}\cos \varphi\right)\vec{e}_{\varphi}
$$

Maintenant, identifions les termes entre parenthèses avec les résultats des dérivées partielles en polaires (relations $\mathbf{(1)}$ et $\mathbf{(3)}$) :

Le coefficient de $\vec{e}_{\rho}$ est $\frac{\partial f}{\partial x}\cos \varphi + \frac{\partial f}{\partial y}\sin \varphi$, qui est exactement $\mathbf{(1)}$ :
    $$
    \frac{\partial f}{\partial \rho}
    $$
Le coefficient de $\vec{e}_{\varphi}$ est $-\frac{\partial f}{\partial x}\sin \varphi + \frac{\partial f}{\partial y}\cos \varphi$, qui est exactement $\mathbf{(3)}$ :
    $$
    \frac{1}{\rho}\frac{\partial f}{\partial \varphi}
    $$


En substituant ces identités, nous obtenons l'expression du gradient en coordonnées polaires :

$$
\vec{\nabla}f = \frac{\partial f}{\partial\rho}\vec{e}_{\rho}+\frac{1}{\rho}\frac{\partial f}{\partial\varphi}\vec{e}_{\varphi}
$$

\subsubsection{Gradient de $f(\rho)=k\rho^{2}/2$}

Le champ $f$ ne dépend que de $\rho$ (la distance à l'origine), donc $\frac{\partial f}{\partial \varphi} = 0$.
$$
f(\rho) = \frac{k}{2}\rho^{2}
$$
Calculons la dérivée par rapport à $\rho$:
$$
\frac{\partial f}{\partial \rho} = \frac{d}{d\rho}\left(\frac{k}{2}\rho^{2}\right) = k\rho
$$
En utilisant l'expression du gradient en coordonnées polaires:
$$
\vec{\nabla}f = \frac{\partial f}{\partial\rho}\vec{e}_{\rho}+\frac{1}{\rho}\frac{\partial f}{\partial\varphi}\vec{e}_{\varphi} = (k\rho)\vec{e}_{\rho} + 0\vec{e}_{\varphi} = k\rho\vec{e}_{\rho}
$$
\section{Plan Tangent}
\subsection{Plan Tangent à $z=f(x,y) = x^2 + xy + y^2$ au point $(1, 2, 7)$}

L'équation du plan tangent à la surface $z=f(x,y)$ au point $(x_0, y_0, z_0)$ est donnée par :
$$
z - z_0 = \frac{\partial f}{\partial x}(x_0, y_0)(x - x_0) + \frac{\partial f}{\partial y}(x_0, y_0)(y - y_0)
$$

A . Calcul des dérivées partielles
$f(x,y) = x^{2}+x y+y^{2}$.
1.  Dérivée par rapport à $x$ :
    $$
    \frac{\partial f}{\partial x} = 2x + y
    $$
2.  Dérivée par rapport à $y$ :
    $$
    \frac{\partial f}{\partial y} = x + 2y
    $$

 B. Évaluation au point $M_0 = (1, 2, 7)$

Le point donné est $M_0=(x_0,y_0,z_0)=(1,2,7)$.

1.  Pente dans la direction $x$ :
    $$
    \frac{\partial f}{\partial x}(1, 2) = 2(1) + 2 = 4
    $$
2.  Pente dans la direction $y$ :
    $$
    \frac{\partial f}{\partial y}(1, 2) = 1 + 2(2) = 5
    $$

3.  Vérification du point : $f(1, 2) = 1^2 + (1)(2) + 2^2 = 1 + 2 + 4 = 7 = z_0$. Le point est bien sur la surface.

 C. Équation du plan tangent

En substituant les valeurs dans l'équation du plan :
$$
z - 7 = 4(x - 1) + 5(y - 2)
$$
Développement et simplification :
$$
z - 7 = 4x - 4 + 5y - 10
$$
$$
4x + 5y - z = 4 + 10 - 7
$$
L'équation du plan tangent est :
$$
\mathbf{4x + 5y - z = 7}
$$

\subsection{Plan Tangent à $z=f(x,y) = 8-x^2-y^2$ au point $(1, 2, 3)$}

L'approche est la même que précédemment, avec $f(x,y)=8-x^{2}-y^{2}$.

A. Calcul des dérivées partielles

1.  Dérivée par rapport à $x$ :
    $$
    \frac{\partial f}{\partial x} = -2x
    $$
2.  Dérivée par rapport à $y$ :
    $$
    \frac{\partial f}{\partial y} = -2y
    $$

B. Évaluation au point $M_0 = (1, 2, 3)$

Le point est $M_0=(x_0,y_0,z_0)=(1,2,3)$.
1.  Pente dans la direction $x$ :
    $$
    \frac{\partial f}{\partial x}(1, 2) = -2(1) = -2
    $$
2.  Pente dans la direction $y$ :
    $$
    \frac{\partial f}{\partial y}(1, 2) = -2(2) = -4
    $$

C. Équation du plan tangent
$$
z - 3 = -2(x - 1) + (-4)(y - 2)
$$
Développement et simplification :
$$
z - 3 = -2x + 2 - 4y + 8
$$
$$
2x + 4y + z = 2 + 8 + 3 = 13
$$
\subsection{Plan Tangent à la surface implicite $f(x,y,z)=0$ au point $M_0=(1, 1, 1)$}

Pour une surface définie implicitement par $f(x,y,z)=C$ (ici $C=0$), le vecteur normal au plan tangent au point $M_0$ est donné par le gradient de $f$ évalué en ce point : $\vec{n} = \vec{\nabla}f(M_0)$.

L'équation du plan tangent est donnée par :
$$
\vec{\nabla}f(M_0) \cdot (\vec{r} - \vec{r}_0) = 0
$$
$$
\frac{\partial f}{\partial x}(M_0)(x - x_0) + \frac{\partial f}{\partial y}(M_0)(y - y_0) + \frac{\partial f}{\partial z}(M_0)(z - z_0) = 0
$$

La fonction est $f(x,y,z)=2x^{3}-3x^{2}+y^{2}z+yz^{2}+\frac{1}{3}z^{3}-\frac{4}{3}$.

 A. Calcul des dérivées partielles

1.  Dérivée par rapport à $x$ :
    $$
    \frac{\partial f}{\partial x} = 6x^2 - 6x
    $$
2.  Dérivée par rapport à $y$ :
    $$
    \frac{\partial f}{\partial y} = 2yz + z^2
    $$
3.  Dérivée par rapport à $z$ :
    $$
    \frac{\partial f}{\partial z} = y^2 + 2yz + z^2
    $$

 B. Évaluation au point $M_0 = (1, 1, 1)$

1.  Composante $x$ du gradient :
    $$
    \frac{\partial f}{\partial x}(1, 1, 1) = 6(1)^2 - 6(1) = 0
    $$
2.  Composante $y$ du gradient :
    $$
    \frac{\partial f}{\partial y}(1, 1, 1) = 2(1)(1) + 1^2 = 3
    $$
3.  Composante $z$ du gradient :
    $$
    \frac{\partial f}{\partial z}(1, 1, 1) = 1^2 + 2(1)(1) + 1^2 = 4
    $$
Le vecteur normal est $\vec{\nabla}f(1, 1, 1) = (0, 3, 4)$.

 C. Équation du plan tangent

Avec $x_0=1, y_0=1, z_0=1$ :
$$
0(x - 1) + 3(y - 1) + 4(z - 1) = 0
$$
$$
3y - 3 + 4z - 4 = 0
$$
L'équation du plan tangent est :
$$
\mathbf{3y + 4z = 7}
$$

\subsection{Plan Tangent à la surface implicite $f(x,y,z)=9$ au point $M_0=(1, -1, 2)$}

La surface est définie par $f(x,y,z)=(x-2)^{2}+(y+3)^{2}+z^{2} = 9$.

L'équation est de la forme $(x-a)^2 + (y-b)^2 + (z-c)^2 = R^2$. Cette équation représente une sphère  de Centre $C = (2, -3, 0)$ et de Rayon $R = \sqrt{9} = 3$


Nous utilisons la méthode du gradient $\vec{\nabla}f$ pour trouver le vecteur normal.

La fonction est $f(x,y,z)=(x-2)^{2}+(y+3)^{2}+z^{2}$.

A. Calcul des dérivées partielles

1.  Dérivée par rapport à $x$ :
    $$
    \frac{\partial f}{\partial x} = 2(x-2)
    $$
2.  Dérivée par rapport à $y$ :
    $$
    \frac{\partial f}{\partial y} = 2(y+3)
    $$
3.  Dérivée par rapport à $z$ :
    $$
    \frac{\partial f}{\partial z} = 2z
    $$

B.  Évaluation au point $M_0 = (1, -1, 2)$

1.  Composante $x$ du gradient :
    $$
    \frac{\partial f}{\partial x}(1, -1, 2) = 2(1 - 2) = -2
    $$
2.  Composante $y$ du gradient :
    $$
    \frac{\partial f}{\partial y}(1, -1, 2) = 2(-1 + 3) = 4
    $$
3.  Composante $z$ du gradient :
    $$
    \frac{\partial f}{\partial z}(1, -1, 2) = 2(2) = 4
    $$
Le vecteur normal est $\vec{n} = (-2, 4, 4)$.

C.  Équation du plan tangent

Avec $x_0=1, y_0=-1, z_0=2$ :
\begin{flalign*}
-2(x - 1) + 4(y - (-1)) + 4(z - 2) &= 0\\
-2x + 2 + 4y + 4 + 4z - 8 &= 0\\
-2x + 4y + 4z &= 8 - 2 - 4\\
-2x + 4y + 4z &= 2\\
x - 2y - 2z &= -1
\end{flalign*}
\section{Champ de gradient}

\subsection{Champ de Vecteurs $\vec{F}(x,y)=(y-3x^2)\vec{e}_{x}+(x+\sin(y))\vec{e}_{y}$}

A. Vérification qu'il est Conservatif
Un champ de vecteurs $\vec{F}(x,y) = P(x,y)\vec{e}_{x} + Q(x,y)\vec{e}_{y}$ défini sur le domaine simplement connexe $\mathbb{R}^2$ est conservatif (ou est un champ de gradient) si son rotationnel est nul, c'est-à-dire si la condition de Schwarz est vérifiée :
$$
\frac{\partial Q}{\partial x} = \frac{\partial P}{\partial y}
$$

Ici, $P(x,y) = y-3x^2$ et $Q(x,y) = x+\sin(y)$.

1.  Calcul de $\frac{\partial P}{\partial y}$ :
    $$
    \frac{\partial P}{\partial y} = \frac{\partial}{\partial y}(y-3x^2) = 1
    $$
2.  Calcul de $\frac{\partial Q}{\partial x}$ :
    $$
    \frac{\partial Q}{\partial x} = \frac{\partial}{\partial x}(x+\sin(y)) = 1
    $$

Puisque $\frac{\partial Q}{\partial x} = \frac{\partial P}{\partial y} = 1$, le champ $\vec{F}$ est bien conservatif sur $\mathbb{R}^2$.

B. Détermination du Potentiel $f$ tel que $\vec{F} = \vec{\nabla}f$

Nous cherchons une fonction scalaire $f(x,y)$ telle que :
$$
\frac{\partial f}{\partial x} = P(x,y) = y-3x^2 \quad \mathbf{(1)}
$$
$$
\frac{\partial f}{\partial y} = Q(x,y) = x+\sin(y) \quad \mathbf{(2)}
$$

1.  Intégration de (1) par rapport à $x$ :
    $$
    f(x,y) = \int (y-3x^2) dx = yx - x^3 + C_1(y)
    $$
    où $C_1(y)$ est une fonction arbitraire dépendant uniquement de $y$ (la "constante" d'intégration).

2.  Dérivation par rapport à $y$ et identification avec (2) :
    Nous dérivons le résultat obtenu par rapport à $y$ :
    $$
    \frac{\partial f}{\partial y} = \frac{\partial}{\partial y}(yx - x^3 + C_1(y)) = x + C_1'(y)
    $$
    En comparant avec l'équation $\mathbf{(2)}$ ($\frac{\partial f}{\partial y} = x+\sin(y)$) :
    $$
    x + C_1'(y) = x + \sin(y)
    $$
    $$
    C_1'(y) = \sin(y)
    $$

3.  Intégration de $C_1'(y)$ pour trouver $C_1(y)$ :
    $$
    C_1(y) = \int \sin(y) dy = -\cos(y) + C
    $$
    où $C$ est une constante réelle.

4.  Conclusion :
    En substituant $C_1(y)$ dans l'expression de $f(x,y)$ :
    $$
    f(x,y) = yx - x^3 - \cos(y) + C
    $$
    Le potentiel scalaire est :
    $$
    \mathbf{f(x,y) = xy - x^3 - \cos(y) + C}
    $$

\subsection{2. Champ de Vecteurs $\vec{F}(x,y,z)=y^2\vec{e}_{x}+(2xy+z)\vec{e}_{y}+(y+3)\vec{e}_{z}$}

A. Vérification qu'il est Conservatif

Un champ de vecteurs $\vec{F}(x,y,z) = P\vec{e}_{x} + Q\vec{e}_{y} + R\vec{e}_{z}$ est conservatif sur $\mathbb{R}^3$ si son rotationnel est nul : $\vec{\nabla} \wedge \vec{F} = \vec{0}$.
Les trois conditions à vérifier sont :
$$
\frac{\partial R}{\partial y} = \frac{\partial Q}{\partial z}, \quad \frac{\partial P}{\partial z} = \frac{\partial R}{\partial x}, \quad \frac{\partial Q}{\partial x} = \frac{\partial P}{\partial y}
$$
Ici, $P = y^2$, $Q = 2xy+z$, $R = y+3$.

1.  Vérification de $\frac{\partial R}{\partial y} = \frac{\partial Q}{\partial z}$ :
     $\frac{\partial R}{\partial y} = \frac{\partial}{\partial y}(y+3) = 1$ et $\frac{\partial Q}{\partial z} = \frac{\partial}{\partial z}(2xy+z) = 1$

2.  Vérification de $\frac{\partial P}{\partial z} = \frac{\partial R}{\partial x}$ :
    $\frac{\partial P}{\partial z} = \frac{\partial}{\partial z}(y^2) = 0$ et $\frac{\partial R}{\partial x} = \frac{\partial}{\partial x}(y+3) = 0$

3.  Vérification de $\frac{\partial Q}{\partial x} = \frac{\partial P}{\partial y}$ :
    $\frac{\partial Q}{\partial x} = \frac{\partial}{\partial x}(2xy+z) = 2y$ et $\frac{\partial P}{\partial y} = \frac{\partial}{\partial y}(y^2) = 2y$

Le champ $\vec{F}$ est bien conservatif sur $\mathbb{R}^3$.

B. Détermination du Potentiel $f$ tel que $\vec{F} = \vec{\nabla}f$

Nous cherchons $f(x,y,z)$ telle que :
$$
\frac{\partial f}{\partial x} = y^2 \quad \mathbf{(1)}
$$
$$
\frac{\partial f}{\partial y} = 2xy+z \quad \mathbf{(2)}
$$
$$
\frac{\partial f}{\partial z} = y+3 \quad \mathbf{(3)}
$$

1.  Intégration de (1) par rapport à $x$ :
    $$
    f(x,y,z) = \int y^2 dx = y^2x + C_1(y,z)
    $$

2.  Dérivation par rapport à $y$ et identification avec (2) :
    $$
    \frac{\partial f}{\partial y} = \frac{\partial}{\partial y}(y^2x + C_1(y,z)) = 2yx + \frac{\partial C_1}{\partial y}
    $$
    En comparant avec $\mathbf{(2)}$ ($\frac{\partial f}{\partial y} = 2xy+z$) :
    $$
    2yx + \frac{\partial C_1}{\partial y} = 2xy+z \quad \Rightarrow \quad \frac{\partial C_1}{\partial y} = z
    $$

3.  Intégration de $\frac{\partial C_1}{\partial y}$ par rapport à $y$ :
    $$
    C_1(y,z) = \int z dy = zy + C_2(z)
    $$
    où $C_2(z)$ est une fonction arbitraire dépendant uniquement de $z$.
    Donc, $f(x,y,z) = y^2x + zy + C_2(z)$.

4.  Dérivation par rapport à $z$ et identification avec (3) :
    $$
    \frac{\partial f}{\partial z} = \frac{\partial}{\partial z}(y^2x + zy + C_2(z)) = y + C_2'(z)
    $$
    En comparant avec $\mathbf{(3)}$ ($\frac{\partial f}{\partial z} = y+3$) :
    $$
    y + C_2'(z) = y + 3 \quad \Rightarrow \quad C_2'(z) = 3
    $$

5.  Intégration de $C_2'(z)$ pour trouver $C_2(z)$ :
    $$
    C_2(z) = \int 3 dz = 3z + C
    $$
    où $C$ est une constante réelle.

6.  Conclusion :
    En substituant $C_2(z)$ dans l'expression de $f(x,y,z)$ :
    $$
    f(x,y,z) = xy^2 + yz + 3z + C
    $$
    Le potentiel scalaire est :
    $$
    \mathbf{f(x,y,z) = xy^2 + yz + 3z + C}
    $$

\subsection{Champ de Vecteurs $\vec{F}(x,y)=\cos x \cos y \vec{e}_{x}-\sin x \sin y \vec{e}_{y}$}

Le domaine est $\mathcal{D}=\{(x,y)\in\mathbb{R}^{2}|-\pi<x<\pi, -\pi/2<y<\pi/2\}$.

 A. Détermination de la Direction Générale de $\vec{F}$
Nous étudions le signe des composantes $P = \cos x \cos y$ et $Q = -\sin x \sin y$ dans chaque case du quadrillage de $\mathcal{D}$ par multiples entiers de $\pi/2$.

!TODO! : Faire le schéma + quadrillage

Le champ est principalement dirigé vers l'Est ($\mathbf{P>0}$) dans la bande $-\pi/2 < x < \pi/2$, et principalement dirigé vers l'Ouest ($\mathbf{P<0}$) en dehors de cette bande.

B. Vérification qu'il est Conservatif

Nous vérifions la condition de Schwarz sur le domaine $\mathcal{D}$ (qui est simplement connexe).
$P(x,y) = \cos x \cos y$ et $Q(x,y) = -\sin x \sin y$.

1.  Calcul de $\frac{\partial P}{\partial y}$ :
    $$
    \frac{\partial P}{\partial y} = \frac{\partial}{\partial y}(\cos x \cos y) = \cos x (-\sin y) = -\cos x \sin y
    $$
2.  Calcul de $\frac{\partial Q}{\partial x}$ :
    $$
    \frac{\partial Q}{\partial x} = \frac{\partial}{\partial x}(-\sin x \sin y) = -(\cos x) \sin y = -\cos x \sin y
    $$

Puisque $\frac{\partial Q}{\partial x} = \frac{\partial P}{\partial y}$, le champ $\vec{F}$ est bien conservatif sur $\mathcal{D}$.

C. Détermination du Potentiel $f$ tel que $\vec{F} = \vec{\nabla}f$

1.  Intégration de $\frac{\partial f}{\partial x} = P$ par rapport à $x$ :
    $$
    f(x,y) = \int \cos x \cos y dx = (\sin x \cos y) + C_1(y)
    $$

2.  Dérivation par rapport à $y$ et identification avec $\frac{\partial f}{\partial y} = Q$ :
    $$
    \frac{\partial f}{\partial y} = \frac{\partial}{\partial y}(\sin x \cos y + C_1(y)) = \sin x (-\sin y) + C_1'(y) = -\sin x \sin y + C_1'(y)
    $$
    En comparant avec $Q = -\sin x \sin y$ :
    $$
    -\sin x \sin y + C_1'(y) = -\sin x \sin y \quad \Rightarrow \quad C_1'(y) = 0
    $$

3.  Intégration de $C_1'(y)$ :
    $$
    C_1(y) = C
    $$

4.  Conclusion :
    $$
    \mathbf{f(x,y) = \sin x \cos y + C}
    $$

D. Détermination des Extrema de $f$ dans $\mathcal{D}$
Les extrema locaux de $f$ sont trouvés lorsque $\vec{\nabla}f = \vec{F} = \vec{0}$.
$$
\begin{cases}
P = \cos x \cos y = 0 \\
Q = -\sin x \sin y = 0
\end{cases}
$$
Dans le domaine $\mathcal{D}$: $-\pi < x < \pi$ et $-\pi/2 < y < \pi/2$.

De $P=0$:
$\cos x = 0 \quad \Rightarrow \quad x = \pm \pi/2$
ou
$\cos y = 0 \quad \Rightarrow \quad y = \pm \pi/2$ (Non inclus dans $\mathcal{D}$)

De $Q=0$:
$\sin x = 0 \quad \Rightarrow \quad x = 0$
ou
$\sin y = 0 \quad \Rightarrow \quad y = 0$

Nous cherchons les points où les deux conditions sont satisfaites.
Si $x = \pm \pi/2$, alors $\sin x = \pm 1$. Pour que $Q=0$, il faut $\sin y = 0$, donc $y=0$.
    Les points critiques sont $\mathbf{( \pi/2, 0 )}$ et $\mathbf{( -\pi/2, 0 )}$.

Évaluons $f(x,y) = \sin x \cos y$ à ces points :

$f(\pi/2, 0) = \sin(\pi/2)\cos(0) = (1)(1) = 1$

$f(-\pi/2, 0) = \sin(-\pi/2)\cos(0) = (-1)(1) = -1$

Analyse des extrema :

Le point $(\pi/2, 0)$ est un Maximum local ($f=1$).

Le point $(-\pi/2, 0)$ est un Minimum local ($f=-1$).


\subsection{Champ de Vecteurs $\vec{F}(x,y)=-y\vec{e}_{x}+x\vec{e}_{y}$}

A. Direction de $\vec{F}$ et Dessin
Le champ de vecteurs est $\vec{F}(x,y) = (-y, x)$.

1.  Relation au vecteur position : Si $\vec{r} = (x, y)$, alors $\vec{F}$ est orthogonal à $\vec{r}$.
    $$
    \vec{F} \cdot \vec{r} = (-y)x + x(y) = -xy + xy = 0
    $$
2.  Direction : Le vecteur $\vec{F}$ est toujours tangent au cercle de rayon $r = \sqrt{x^2+y^2}$ passant par le point $(x,y)$. Le champ est dirigé dans le sens trigonométrique.



B. $\vec{F}$ est-il Conservatif ?

Nous vérifions la condition de Schwarz.
$P(x,y) = -y$ et $Q(x,y) = x$.

1.  Calcul de $\frac{\partial P}{\partial y}$ :
    $$
    \frac{\partial P}{\partial y} = \frac{\partial}{\partial y}(-y) = -1
    $$
2.  Calcul de $\frac{\partial Q}{\partial x}$ :
    $$
    \frac{\partial Q}{\partial x} = \frac{\partial}{\partial x}(x) = 1
    $$

Puisque $\frac{\partial Q}{\partial x} \ne \frac{\partial P}{\partial y}$ ($1 \ne -1$), le champ $\vec{F}$ n'est pas conservatif.

Un champ conservatif n'a pas de circulation le long d'une courbe fermée. Or, ce champ $\vec{F}$ est un champ de rotation pure ; le travail effectué le long de tout cercle centré à l'origine est non nul. Cela confirme qu'il n'est pas conservatif.


\subsection{Champ de Vecteurs $\vec{F}(x,y)=-\frac{y}{x^{2}+y^{2}}\vec{e}_{x}+\frac{x}{x^{2}+y^{2}}\vec{e}_{y}$}


A. Est-il Conservatif sur $\mathcal{D}$ ?

Nous vérifions la condition de Schwarz sur $\mathcal{D}$.
$P(x,y) = -\frac{y}{x^{2}+y^{2}}$ et $Q(x,y) = \frac{x}{x^{2}+y^{2}}$.

1.  Calcul de $\frac{\partial P}{\partial y}$ :
    $$
    \frac{\partial P}{\partial y} = -\left[\frac{(1)(x^2+y^2) - y(2y)}{(x^2+y^2)^2}\right] = -\frac{x^2+y^2-2y^2}{(x^2+y^2)^2} = -\frac{x^2-y^2}{(x^2+y^2)^2} = \frac{y^2-x^2}{(x^2+y^2)^2}
    $$
2.  Calcul de $\frac{\partial Q}{\partial x}$ :
    $$
    \frac{\partial Q}{\partial x} = \frac{(1)(x^2+y^2) - x(2x)}{(x^2+y^2)^2} = \frac{x^2+y^2-2x^2}{(x^2+y^2)^2} = \frac{y^2-x^2}{(x^2+y^2)^2}
    $$

Puisque $\frac{\partial Q}{\partial x} = \frac{\partial P}{\partial y}$, le champ $\vec{F}$ satisfait la condition de rotationnel nul partout où il est défini (sur $\mathcal{D}$).

Comme le domaine $\mathcal{D}$  n'est pas simplement connexe, le fait que le rotationnel soit nul ne garantit pas que le champ soit conservatif.

B. Détermination de $f$ tel que $\vec{F} = \vec{\nabla}f$
En coordonnées polaires, $\vec{F}$ s'exprime par :
$$
\vec{F} = \frac{1}{\rho}\vec{e}_{\varphi}
$$
(où $\frac{1}{\rho} \vec{e}_{\varphi}$ est la composante orthoradiale).

Si $\vec{F}$ dérivait d'un potentiel $f(\rho, \varphi)$, nous aurions :
$$
\vec{F} = \frac{\partial f}{\partial\rho}\vec{e}_{\rho}+\frac{1}{\rho}\frac{\partial f}{\partial\varphi}\vec{e}_{\varphi}
$$
Par identification :
$$
\frac{\partial f}{\partial\rho} = 0 \quad \mathbf{(1')}
$$
$$
\frac{1}{\rho}\frac{\partial f}{\partial\varphi} = \frac{1}{\rho} \quad \Rightarrow \quad \frac{\partial f}{\partial\varphi} = 1 \quad \mathbf{(2')}
$$

1.  Intégration de (1') : $f$ ne dépend pas de $\rho$ :
    $$
    f(x,y) = C_1(\varphi)
    $$
2.  Intégration de (2') :
    $$
    f(\varphi) = \int 1 d\varphi = \varphi + C
    $$
    où $\varphi = \arctan(y/x)$ est l'angle en coordonnées polaires.

Le potentiel est :
$$
f(x,y) = \arctan(y/x) + C = \varphi + C
$$
Le potentiel est l'angle $\varphi$. Cependant, l'angle $\varphi$ n'est pas une fonction uniforme sur le domaine $\mathcal{D}$, car il change de $2\pi$ après un tour complet autour de l'origine. L'incompatibilité entre $\vec{\nabla}f$ (qui a un rotationnel nul) et l'intégrale de circulation non nulle est due à cette singularité de $\varphi$ à l'origine, ce qui explique pourquoi un potentiel $f$ n'existe pas sur l'ensemble du domaine $\mathcal{D}$.
\section{Calculs de divergences et de rotationnels}
\subsection{Champ de Vecteurs 1}

Le champ $\vec{F}$ est défini par $\vec{F} = (x^2+y^2)^a \vec{r}_{\perp}$, où $\vec{r}_{\perp} = x\vec{e}_{x}+y\vec{e}_{y}$ est la composante radiale du vecteur position dans le plan $(O, \vec{e}_x, \vec{e}_y)$. En coordonnées polaires, $x^2+y^2 = \rho^2$, et $\vec{r}_{\perp} = \rho \vec{e}_{\rho}$.

$$
\vec{F} = (\rho^2)^a (\rho \vec{e}_{\rho}) = \rho^{2a+1} \vec{e}_{\rho}
$$
Ce champ est  radial.

A. Description du Champ, Divergence et Rotationnel

$\vec{F}$ est un champ radial centré sur l'axe $(O, \vec{e}_z)$. Sa direction est toujours le long du rayon $\vec{e}_{\rho}$ (s'éloignant de l'axe $z$). Sa magnitude, $\rho^{2a+1}$, ne dépend que de la distance $\rho$ à l'axe $z$.

Calcul de la Divergence :

    En coordonnées cartésiennes :
    $$
    \operatorname{div}\vec{F} = \frac{\partial F_x}{\partial x} + \frac{\partial F_y}{\partial y} + \frac{\partial F_z}{\partial z}
    $$
    $F_x = x(x^2+y^2)^a$, $F_y = y(x^2+y^2)^a$, $F_z = 0$.
    $$
    \frac{\partial F_x}{\partial x} = (x^2+y^2)^a + x \cdot a(x^2+y^2)^{a-1}(2x) = (x^2+y^2)^a + 2ax^2(x^2+y^2)^{a-1}
    $$
    $$
    \frac{\partial F_y}{\partial y} = (x^2+y^2)^a + y \cdot a(x^2+y^2)^{a-1}(2y) = (x^2+y^2)^a + 2ay^2(x^2+y^2)^{a-1}
    $$
    $$
    \operatorname{div}\vec{F} = 2(x^2+y^2)^a + 2a(x^2+y^2)^{a-1}(x^2+y^2) = = 2(1+a)(x^2+y^2)^a
    $$

Calcul du Rotationnel :

    Comme $\vec{F}$ est purement radial ($\vec{F}$ est colinéaire à $\vec{r}_{\perp}$) et ne dépend que de $\rho$, son rotationnel est nul.
    $$
    \mathbf{\operatorname{rot}\vec{F} = \vec{0}}
    $$

    En effet, $\vec{F}$ est un champ de gradient : $\vec{F} = \vec{\nabla}f$ avec $f(\rho) = \frac{\rho^{2a+2}}{2a+2}$ (pour $a \ne -1$).


\subsection{Champ de Vecteurs2}

Le champ $\vec{F}$ est défini par $\vec{F} = (x^2+y^2)^a \vec{v}_{rot}$, où $\vec{v}_{rot} = -y\vec{e}_{x}+x\vec{e}_{y}$ est le champ de rotation. En coordonnées polaires, $\vec{v}_{rot} = \rho \vec{e}_{\varphi}$.

$$
\vec{F} = (\rho^2)^a (\rho \vec{e}_{\varphi}) = \rho^{2a+1} \vec{e}_{\varphi}
$$
Ce champ est purement orthoradial.

A. Description du Champ, Divergence et Rotationnel

$\vec{F}$ est un champ de rotation pure (circulaire) autour de l'axe $(O, \vec{e}_z)$. Sa direction est toujours tangentielle ($\vec{e}_{\varphi}$), dans le sens anti-horaire. Sa magnitude, $\rho^{2a+1}$, ne dépend que de la distance $\rho$ à l'axe $z$.

Calcul de la Divergence :

    En utilisant l'expression en coordonnées cylindriques (car $A=0$ et $C=0$) :
    $A = 0$, $B = \rho^{2a+1}$, $C = 0$.
    $$
    \operatorname{div}\vec{F} = \frac{1}{\rho}\frac{\partial}{\partial\rho}(\rho A) + \frac{1}{\rho}\frac{\partial B}{\partial\varphi} + \frac{\partial C}{\partial z}
    $$
    Puisque $B$ ne dépend que de $\rho$ (et non de $\varphi$), $\frac{\partial B}{\partial\varphi} = 0$.
    $$
    \operatorname{div}\vec{F} = 0 + 0 + 0 = 0
    $$
    Interprétation : Le champ est incompressible (pas de sources ni de puits).
    $$
    \mathbf{\operatorname{div}\vec{F} = 0}
    $$

Calcul du Rotationnel :

    En coordonnées cartésiennes, seule la composante $z$ est non nulle.
    $$
    (\operatorname{rot}\vec{F})_z = \frac{\partial Q}{\partial x} - \frac{\partial P}{\partial y}
    $$
    $P = -y(x^2+y^2)^a$, $Q = x(x^2+y^2)^a$.

    $$
    \frac{\partial Q}{\partial x} = (x^2+y^2)^a + x \cdot a(x^2+y^2)^{a-1}(2x) = (x^2+y^2)^a + 2ax^2(x^2+y^2)^{a-1}
    $$
    $$
    \frac{\partial P}{\partial y} = -(x^2+y^2)^a - y \cdot a(x^2+y^2)^{a-1}(2y) = -(x^2+y^2)^a - 2ay^2(x^2+y^2)^{a-1}
    $$
    $$
    (\operatorname{rot}\vec{F})_z = \left( (x^2+y^2)^a + 2ax^2(x^2+y^2)^{a-1} \right) - \left( -(x^2+y^2)^a - 2ay^2(x^2+y^2)^{a-1} \right)
    $$
    $$
    (\operatorname{rot}\vec{F})_z = 2(x^2+y^2)^a + 2a(x^2+y^2)^{a-1}(x^2+y^2)
    $$
    $$
    (\operatorname{rot}\vec{F})_z = 2(x^2+y^2)^a + 2a(x^2+y^2)^a = 2(1+a)(x^2+y^2)^a
    $$
    $$
    \mathbf{\operatorname{rot}\vec{F} = 2(1+a)(x^2+y^2)^a \vec{e}_z}
    $$

    Remarque : Le résultat de $\operatorname{rot}\vec{F}$ pour le champ de rotation (Q2) est égal au résultat de $\operatorname{div}\vec{F}$ pour le champ radial (Q1).

\subsection{Champ de Vecteurs 3}

Le champ est le produit vectoriel entre un vecteur fixé $\vec{u}$ et le vecteur position $\vec{r}=(x,y,z)$.

A. Description du Champ

Le champ $\vec{F}$ est orthogonal à la fois à $\vec{u}$ et à $\vec{r}$. $\vec{F}$ décrit un mouvement de rotation autour d'un axe passant par l'origine et parallèle à $\vec{u}$.

La magnitude de $\vec{F}$ est $|\vec{F}| = |\vec{u}| |\vec{r}| \sin \alpha$, où $\alpha$ est l'angle entre $\vec{u}$ et $\vec{r}$. Si $\rho$ est la distance du point $\vec{r}$ à l'axe de rotation (l'axe $\vec{u}$), alors $\rho = |\vec{r}| \sin \alpha$.

Donc $|\vec{F}| = |\vec{u}| \rho$. La vitesse de rotation est proportionnelle à la distance $\rho$ à l'axe, caractéristique d'un corps rigide en rotation avec une vitesse angulaire $\vec{\omega} = -\vec{u}$.

B. Calcul de la Divergence

En utilisant la formule $ \operatorname{div}(\vec{A} \wedge \vec{B}) = \vec{B} \cdot (\operatorname{rot}\vec{A}) - \vec{A} \cdot (\operatorname{rot}\vec{B})$:

Ici, $\vec{A}=\vec{u}$ (constant) et $\vec{B}=\vec{r}$.

$\operatorname{rot}\vec{u} = \vec{0}$ (rotationnel d'un vecteur constant est nul)

$\operatorname{rot}\vec{r} = \vec{0}$ (rotationnel de $\vec{r}$ est nul)

$$
\operatorname{div}\vec{F} = \operatorname{div}(\vec{u} \wedge \vec{r}) = \vec{r} \cdot (\operatorname{rot}\vec{u}) - \vec{u} \cdot (\operatorname{rot}\vec{r}) = \vec{r} \cdot \vec{0} - \vec{u} \cdot \vec{0} = 0
$$
$$
\mathbf{\operatorname{div}\vec{F} = 0}
$$

C. Calcul du Rotationnel

En utilisant la formule $\operatorname{rot}(\vec{A} \wedge \vec{B}) = \vec{A} (\operatorname{div}\vec{B}) - \vec{B} (\operatorname{div}\vec{A}) + (\vec{B} \cdot \vec{\nabla})\vec{A} - (\vec{A} \cdot \vec{\nabla})\vec{B}$ :

$\operatorname{div}\vec{r} = \frac{\partial x}{\partial x} + \frac{\partial y}{\partial y} + \frac{\partial z}{\partial z} = 1+1+1 = 3$

$\operatorname{div}\vec{u} = 0$

$(\vec{r} \cdot \vec{\nabla})\vec{u} = 0$ (car $\vec{u}$ est constant)

$(\vec{u} \cdot \vec{\nabla})\vec{r} = \vec{u}$ (aussi car $\vec{u}$ est constant)

$$
\operatorname{rot}\vec{F} = \vec{u}(3) - \vec{r}(0) + \vec{0} - \vec{u} = 3\vec{u} - \vec{u} = 2\vec{u}
$$
$$
\mathbf{\operatorname{rot}\vec{F} = 2\vec{u}}
$$
\section{Identités entre rotationnel et divergent}
\subsection{Identité 1}
Montrer par un calcul direct que $\operatorname{rot}(\vec{\nabla}f)=\vec{0}$

Le gradient d'un champ scalaire $f(x,y,z)$ est $\vec{\nabla}f = \left(\frac{\partial f}{\partial x}, \frac{\partial f}{\partial y}, \frac{\partial f}{\partial z}\right)$.

Le rotationnel du gradient est donné par :
$$
\operatorname{rot}(\vec{\nabla}f) = \vec{\nabla} \wedge (\vec{\nabla}f) = \begin{vmatrix} \vec{e}_x & \vec{e}_y & \vec{e}_z \\ \frac{\partial}{\partial x} & \frac{\partial}{\partial y} & \frac{\partial}{\partial z} \\ \frac{\partial f}{\partial x} & \frac{\partial f}{\partial y} & \frac{\partial f}{\partial z} \end{vmatrix}
$$

Calculons la composante en $\vec{e}_x$ :
$$
(\operatorname{rot}(\vec{\nabla}f))_x = \frac{\partial}{\partial y}\left(\frac{\partial f}{\partial z}\right) - \frac{\partial}{\partial z}\left(\frac{\partial f}{\partial y}\right) = \frac{\partial^{2}f}{\partial y \partial z} - \frac{\partial^{2}f}{\partial z \partial y}
$$
En supposant que $f$ est suffisamment régulière (Théorème de Schwarz), les dérivées croisées sont égales : $\frac{\partial^{2}f}{\partial y \partial z} = \frac{\partial^{2}f}{\partial z \partial y}$.

D'où :
$$
(\operatorname{rot}(\vec{\nabla}f))_x = 0
$$
Les composantes en $\vec{e}_y$ et $\vec{e}_z$ s'annulent de la même manière.
$$
\mathbf{\operatorname{rot}(\vec{\nabla}f) = \vec{0}}
$$
Interprétation : Un champ de vecteurs dérivant d'un potentiel scalaire (champ conservatif) n'a pas de rotation.

\subsection{Identité 2}
Montrer par un calcul direct que $\operatorname{div}(\operatorname{rot}\vec{F})=0$

Nous calculons la divergence du rotationnel d'un champ vectoriel $\vec{F} = F_x \vec{e}_x + F_y \vec{e}_y + F_z \vec{e}_z$.
$$
\operatorname{div}(\operatorname{rot}\vec{F}) = \frac{\partial}{\partial x}(\operatorname{rot}\vec{F})_x + \frac{\partial}{\partial y}(\operatorname{rot}\vec{F})_y + \frac{\partial}{\partial z}(\operatorname{rot}\vec{F})_z
$$
$$
\operatorname{div}(\operatorname{rot}\vec{F}) = \frac{\partial}{\partial x}\left(\frac{\partial F_z}{\partial y} - \frac{\partial F_y}{\partial z}\right) + \frac{\partial}{\partial y}\left(\frac{\partial F_x}{\partial z} - \frac{\partial F_z}{\partial x}\right) + \frac{\partial}{\partial z}\left(\frac{\partial F_y}{\partial x} - \frac{\partial F_x}{\partial y}\right)
$$
En développant et en réordonnant les termes :
$$
\operatorname{div}(\operatorname{rot}\vec{F}) = \left(\frac{\partial^{2}F_z}{\partial x \partial y} - \frac{\partial^{2}F_z}{\partial y \partial x}\right) + \left(\frac{\partial^{2}F_x}{\partial y \partial z} - \frac{\partial^{2}F_x}{\partial z \partial y}\right) + \left(\frac{\partial^{2}F_y}{\partial z \partial x} - \frac{\partial^{2}F_y}{\partial x \partial z}\right)
$$
Par le Théorème de Schwarz, toutes les dérivées croisées s'annulent :
$$
\operatorname{div}(\operatorname{rot}\vec{F}) = 0 + 0 + 0
$$
D'où :
$$
\mathbf{\operatorname{div}(\operatorname{rot}\vec{F}) = 0}
$$
\subsection{Identité 3}
Vérifier par un calcul direct que $\operatorname{div}(\vec{\nabla}f)=\Delta f$

Le Laplacien $\Delta$ est défini comme $\Delta=\vec{\nabla}\cdot\vec{\nabla}=\frac{\partial^{2}}{\partial x^{2}}+\frac{\partial^{2}}{\partial y^{2}}+\frac{\partial^{2}}{\partial z^{2}}$.

Nous calculons la divergence du gradient:
$$
\operatorname{div}(\vec{\nabla}f) = \vec{\nabla} \cdot \left(\frac{\partial f}{\partial x}\vec{e}_{x} + \frac{\partial f}{\partial y}\vec{e}_{y} + \frac{\partial f}{\partial z}\vec{e}_{z}\right)
$$
$$
\operatorname{div}(\vec{\nabla}f) = \frac{\partial}{\partial x}\left(\frac{\partial f}{\partial x}\right) + \frac{\partial}{\partial y}\left(\frac{\partial f}{\partial y}\right) + \frac{\partial}{\partial z}\left(\frac{\partial f}{\partial z}\right)
$$
$$
\operatorname{div}(\vec{\nabla}f) = \frac{\partial^{2}f}{\partial x^{2}} + \frac{\partial^{2}f}{\partial y^{2}} + \frac{\partial^{2}f}{\partial z^{2}}
$$
Ceci est la définition du Laplacien $\Delta f$ :
$$
\mathbf{\operatorname{div}(\vec{\nabla}f) = \Delta f}
$$

\subsection{Identité 4}
Vérifier par un calcul direct que $\operatorname{div}(f\vec{F})=f \operatorname{div}\vec{F}+(\vec{\nabla}f)\cdot\vec{F}$

Nous calculons la divergence du produit du champ scalaire $f$ et du champ vectoriel $\vec{F} = F_x \vec{e}_x + F_y \vec{e}_y + F_z \vec{e}_z$.

$$
\operatorname{div}(f\vec{F}) = \frac{\partial}{\partial x}(fF_x) + \frac{\partial}{\partial y}(fF_y) + \frac{\partial}{\partial z}(fF_z)
$$
En appliquant la règle du produit $\frac{\partial}{\partial x}(uv) = \frac{\partial u}{\partial x}v + u\frac{\partial v}{\partial x}$ à chaque terme :
$$
\operatorname{div}(f\vec{F}) = \left(\frac{\partial f}{\partial x}F_x + f\frac{\partial F_x}{\partial x}\right) + \left(\frac{\partial f}{\partial y}F_y + f\frac{\partial F_y}{\partial y}\right) + \left(\frac{\partial f}{\partial z}F_z + f\frac{\partial F_z}{\partial z}\right)
$$
Regroupons les termes contenant $f$ et ceux contenant les composantes de $\vec{F}$ :
$$
\operatorname{div}(f\vec{F}) = f\left(\frac{\partial F_x}{\partial x} + \frac{\partial F_y}{\partial y} + \frac{\partial F_z}{\partial z}\right) + \left(\frac{\partial f}{\partial x}F_x + \frac{\partial f}{\partial y}F_y + \frac{\partial f}{\partial z}F_z\right)
$$
Par définition :
$$
\left(\frac{\partial F_x}{\partial x} + \frac{\partial F_y}{\partial y} + \frac{\partial F_z}{\partial z}\right) = \operatorname{div}\vec{F}
$$
$$
\left(\frac{\partial f}{\partial x}F_x + \frac{\partial f}{\partial y}F_y + \frac{\partial f}{\partial z}F_z\right) = \vec{\nabla}f \cdot \vec{F}
$$
D'où l'identité :
$$
\mathbf{\operatorname{div}(f\vec{F}) = f \operatorname{div}\vec{F}+(\vec{\nabla}f)\cdot\vec{F}}
$$

\subsection{Identité 5}
Vérifier par un calcul direct que $\operatorname{rot}(\operatorname{rot}\vec{F})=\vec{\nabla}\operatorname{div}\vec{F}-\Delta\vec{F}$

C'est l'identité de Helmholtz. Nous vérifions la composante $x$ (les autres composantes sont vérifiées de manière similaire).

$$
(\operatorname{rot}(\operatorname{rot}\vec{F}))_x = \frac{\partial}{\partial y}(\operatorname{rot}\vec{F})_z - \frac{\partial}{\partial z}(\operatorname{rot}\vec{F})_y
$$
En substituant les composantes du rotationnel :
$$
(\operatorname{rot}(\operatorname{rot}\vec{F}))_x = \frac{\partial}{\partial y}\left(\frac{\partial F_y}{\partial x} - \frac{\partial F_x}{\partial y}\right) - \frac{\partial}{\partial z}\left(\frac{\partial F_x}{\partial z} - \frac{\partial F_z}{\partial x}\right)
$$
$$
(\operatorname{rot}(\operatorname{rot}\vec{F}))_x = \frac{\partial^{2} F_y}{\partial y \partial x} - \frac{\partial^{2} F_x}{\partial y^{2}} - \frac{\partial^{2} F_x}{\partial z^{2}} + \frac{\partial^{2} F_z}{\partial z \partial x}
$$
Réarrangeons les termes et ajoutons/soustraions $\frac{\partial^{2} F_x}{\partial x^{2}}$ :
$$
(\operatorname{rot}(\operatorname{rot}\vec{F}))_x = \left(\frac{\partial^{2} F_x}{\partial x^{2}} + \frac{\partial^{2} F_y}{\partial y \partial x} + \frac{\partial^{2} F_z}{\partial z \partial x}\right) - \left(\frac{\partial^{2} F_x}{\partial x^{2}} + \frac{\partial^{2} F_x}{\partial y^{2}} + \frac{\partial^{2} F_x}{\partial z^{2}}\right)
$$

Le premier terme est la dérivée par rapport à $x$ de la divergence de $\vec{F}$ (par Schwarz, $\partial_{yx}^2 = \partial_{xy}^2$ et $\partial_{zx}^2 = \partial_{xz}^2$) :
$$
\frac{\partial}{\partial x}\left(\frac{\partial F_x}{\partial x} + \frac{\partial F_y}{\partial y} + \frac{\partial F_z}{\partial z}\right) = \frac{\partial}{\partial x}(\operatorname{div}\vec{F}) = (\vec{\nabla}\operatorname{div}\vec{F})_x
$$
Le second terme est le Laplacien appliqué à la composante $x$ de $\vec{F}$ :
$$
\Delta F_x = (\Delta\vec{F})_x
$$
Donc :
$$
(\operatorname{rot}(\operatorname{rot}\vec{F}))_x = (\vec{\nabla}\operatorname{div}\vec{F})_x - (\Delta\vec{F})_x
$$
L'identité est ainsi vérifiée pour la composante $x$ (et par extension pour les autres) :
$$
\mathbf{\operatorname{rot}(\operatorname{rot}\vec{F}) = \vec{\nabla}\operatorname{div}\vec{F} - \Delta\vec{F}}
$$

\subsection{Expression non définie 1}

Soit f un champ scalaire. Que peut-on dire de l'expression $\operatorname{rot}f-\operatorname{div}f$?

L'expression $\operatorname{rot}f - \operatorname{div}f$ n'a pas de sens physique ou mathématique.

1.  Rotationnel : L'opérateur $\operatorname{rot}$ (rotationnel) s'applique uniquement à un champ de vecteurs ($\operatorname{rot}\vec{F}$) et donne un vecteur.
2.  Divergence : L'opérateur $\operatorname{div}$ (divergence) s'applique uniquement à un champ de vecteurs ($\operatorname{div}\vec{F}$) et donne un scalaire.

On ne peut pas effectuer une soustraction entre ces deux termes car ils sont mal définis.
\section{Formule de Green}
\subsection{Champ 1}

Le champ est $\vec{F}(x,y) = P\vec{e}_x + Q\vec{e}_y$ avec $P(x,y) = -y^3$ et $Q(x,y) = 2$.
La courbe $C$ est le cercle $x^2 + y^2 = 4$ (rayon $R=2$)[cite: 78, 79]. $\mathcal{R}$ est le disque $x^2 + y^2 \le 4$.

A. Calcul de la circulation $\int_{\mathcal{C}}\vec{F}\cdot d\vec{r}$

Paramétrisation de $C$ :
$$
x = 2\cos t, \quad y = 2\sin t, \quad 0 \le t \le 2\pi
$$
$$
dx = -2\sin t dt, \quad dy = 2\cos t dt
$$
Intégrale :
$$
\int_{\mathcal{C}}\vec{F}\cdot d\vec{r} = \int_{0}^{2\pi} [P dx + Q dy] = \int_{0}^{2\pi} [(-y^3)(-2\sin t) + (2)(2\cos t)] dt
$$
$$
= \int_{0}^{2\pi} [-(2\sin t)^3 (-2\sin t) + 4\cos t] dt = \int_{0}^{2\pi} [16\sin^4 t + 4\cos t] dt
$$
L'intégrale de $4\cos t$ sur $[0, 2\pi]$ est nulle. Pour $\sin^4 t$, on utilise l'identité $\int_{0}^{2\pi} \sin^4 t dt = \frac{3\pi}{4}$.
$$
\int_{\mathcal{C}}\vec{F}\cdot d\vec{r} = 16 \cdot \frac{3\pi}{4} + 0 = \mathbf{12\pi}
$$

B. Calcul de l'intégrale double $\iint_{\mathcal{R}}(\frac{\partial Q}{\partial x}-\frac{\partial P}{\partial y})dx dy$

Rotationnel 2D :
$$
\frac{\partial Q}{\partial x} = \frac{\partial}{\partial x}(2) = 0
$$
$$
\frac{\partial P}{\partial y} = \frac{\partial}{\partial y}(-y^3) = -3y^2
$$
$$
\frac{\partial Q}{\partial x}-\frac{\partial P}{\partial y} = 0 - (-3y^2) = 3y^2
$$
Intégrale double (en coordonnées polaires) : $x=\rho \cos \varphi$, $y=\rho \sin \varphi$, $dA = \rho d\rho d\varphi$.
$$
\iint_{\mathcal{R}} 3y^2 dx dy = \int_{0}^{2\pi} \int_{0}^{2} 3(\rho \sin \varphi)^2 \rho d\rho d\varphi = 3 \left(\int_{0}^{2} \rho^3 d\rho\right) \left(\int_{0}^{2\pi} \sin^2 \varphi d\varphi\right)
$$
$$
= 3 \cdot \left[\frac{\rho^4}{4}\right]_{0}^{2} \cdot \left[\pi\right]
$$
(où $\int_{0}^{2\pi} \sin^2 \varphi d\varphi = \pi$)
$$
= 3 \cdot \left(\frac{16}{4}\right) \cdot \pi = 3 \cdot 4 \cdot \pi = \mathbf{12\pi}
$$
Conclusion : Ok.

\subsection{Champ 2}
Le champ est $P(x,y) = y^2$ et $Q(x,y) = -3x$. $C$ est le cercle unité $x^2+y^2=1$. La paramétrisation $\vec{r}(t)=\sin(t)\vec{e}_{x}-\cos(t)\vec{e}_{y}$ pour $t \in [-\pi, \pi]$ correspond au parcours du cercle dans le sens trigo.

A. Calcul par intégrale curviligne

Paramétrisation : $x=\sin t$, $y=-\cos t$. $dx=\cos t dt$, $dy=\sin t dt$.
$$
\int_{\mathcal{C}}\vec{F}\cdot d\vec{r} = \int_{-\pi}^{\pi} [P dx + Q dy] = \int_{-\pi}^{\pi} [(-\cos t)^2 \cos t + (-3\sin t) \sin t] dt
$$
$$
= \int_{-\pi}^{\pi} [\cos^3 t - 3\sin^2 t] dt
$$

$\int_{-\pi}^{\pi} \cos^3 t dt = 0$.

$\int_{-\pi}^{\pi} -3\sin^2 t dt = -3 \int_{-\pi}^{\pi} \frac{1-\cos(2t)}{2} dt = -3 \left[\frac{t}{2} - \frac{\sin(2t)}{4}\right]_{-\pi}^{\pi} = -3(\pi) = -3\pi$
$$
\int_{\mathcal{C}}\vec{F}\cdot d\vec{r} = 0 - 3\pi = \mathbf{-3\pi}
$$

B. Vérification par la Formule de Green
Rotationnel 2D :
$$
\frac{\partial Q}{\partial x} = -3, \quad \frac{\partial P}{\partial y} = 2y
$$
$$
\frac{\partial Q}{\partial x}-\frac{\partial P}{\partial y} = -3 - 2y
$$
Intégrale double : L'aire de $\mathcal{R}$ est $\pi$.
$$
\iint_{\mathcal{R}} (-3 - 2y) dx dy = \iint_{\mathcal{R}} -3 dx dy - \iint_{\mathcal{R}} 2y dx dy
$$
$\iint_{\mathcal{R}} -3 dx dy = -3 \cdot \text{Aire}(\mathcal{R}) = -3\pi$

$\iint_{\mathcal{R}} 2y dx dy = 0$ (par symétrie du disque $\mathcal{R}$ et $2y$
est impaire en $y$)

$$
\iint_{\mathcal{R}}\left(\frac{\partial Q}{\partial x} - \frac{\partial P}{\partial y}\right) dx dy = -3\pi + 0 = \mathbf{-3\pi}
$$
Conclusion : OK.
\subsection{Champ 3}

Le champ est $P(x,y) = y-x$ et $Q(x,y) = x-3y$. $C$ est l'ellipse $x=4\cos t, y=3\sin t$ dans le sens positif. $\mathcal{R}$ est la région elliptique.

A. Calcul par intégrale curviligne

Paramétrisation : $x=4\cos t$, $y=3\sin t$. $dx=-4\sin t dt$, $dy=3\cos t dt$.
$$
\int_{\mathcal{C}}\vec{F}\cdot d\vec{r} = \int_{0}^{2\pi} [(y-x) dx + (x-3y) dy]
$$
$$
= \int_{0}^{2\pi} [(3\sin t - 4\cos t)(-4\sin t) + (4\cos t - 9\sin t)(3\cos t)] dt
$$
$$
= \int_{0}^{2\pi} [-12\sin^2 t + 16\sin t \cos t + 12\cos^2 t - 27\sin t \cos t] dt
$$
$$
= \int_{0}^{2\pi} [12(\cos^2 t - \sin^2 t) - 11\sin t \cos t] dt = \int_{0}^{2\pi} [12\cos(2t) - \frac{11}{2}\sin(2t)] dt
$$
L'intégrale des fonctions trigonométriques sur une période entière est nulle.
$$
\int_{\mathcal{C}}\vec{F}\cdot d\vec{r} = 0 - 0 = \mathbf{0}
$$

B. Vérification par la Formule de Green

Rotationnel 2D :
$$
\frac{\partial Q}{\partial x} = \frac{\partial}{\partial x}(x-3y) = 1
$$
$$
\frac{\partial P}{\partial y} = \frac{\partial}{\partial y}(y-x) = 1
$$
$$
\frac{\partial Q}{\partial x}-\frac{\partial P}{\partial y} = 1 - 1 = \mathbf{0}
$$
Intégrale double :
$$
\iint_{\mathcal{R}}\left(\frac{\partial Q}{\partial x} - \frac{\partial P}{\partial y}\right) dx dy = \iint_{\mathcal{R}} 0 dx dy = \mathbf{0}
$$
Conclusion : OK.


\subsection{Champ 4}

Le champ est $P(x,y)=\cos y$ et $Q(x,y)=\sin x$. $C$ est le bord du triangle $R$.
Pour que la formule de Green fonctionne, $C$ doit être parcouru dans le sens positif (région $R$ à gauche). L'ordre positif est $\mathbf{C(0,0) \to A(\pi/4, 0) \to B(0, \pi/4) \to C(0,0)}$.

A. Calcul par intégrale curviligne

$C \to A$ (Segment $y=0$): $x: 0 \to \pi/4$. $dy=0$.
    $$
    \int_{C \to A} \vec{F}\cdot d\vec{r} = \int_{0}^{\pi/4} \cos(0) dx = \int_{0}^{\pi/4} 1 dx = \pi/4
    $$
$A \to B$ (Segment $x+y=\pi/4$): $x: \pi/4 \to 0$. $y=\pi/4-x$, $dy=-dx$.
    $$
    \int_{A \to B} \vec{F}\cdot d\vec{r} = \int_{\pi/4}^{0} [\cos(\pi/4-x) - \sin x] dx = -\left( \sqrt{2}-1 \right) = 1 - \sqrt{2}
    $$
$B \to C$ (Segment $x=0$): $y: \pi/4 \to 0$. $dx=0$.
    $$
    \int_{B \to C} \vec{F}\cdot d\vec{r} = \int_{\pi/4}^{0} \sin(0) dy = 0
    $$
$$
\int_{\mathcal{C}}\vec{F}\cdot d\vec{r} = \pi/4 + (1 - \sqrt{2}) + 0 = \mathbf{1 + \pi/4 - \sqrt{2}}
$$

B.  Vérification par la Formule de Green

Rotationnel 2D :
$$
\frac{\partial Q}{\partial x} = \cos x, \quad \frac{\partial P}{\partial y} = -\sin y
$$
$$
\frac{\partial Q}{\partial x}-\frac{\partial P}{\partial y} = \cos x + \sin y
$$
Intégrale double : $\mathcal{R}: 0 \le x \le \pi/4$ et $0 \le y \le \pi/4 - x$.
$$
\iint_{\mathcal{R}} (\cos x + \sin y) dx dy = \int_{0}^{\pi/4} \left( \int_{0}^{\pi/4-x} (\cos x + \sin y) dy \right) dx
$$
Après intégration, on trouve :
$$
\iint_{\mathcal{R}} (\cos x + \sin y) dx dy = \mathbf{1 + \pi/4 - \sqrt{2}}
$$
Conclusion : OK.

\subsection{Champ 5}

Le champ est $P(x,y)=y$ et $Q(x,y)=-x/2$. $C$ est le contour de la région $\mathcal{R}$ définie par $0 \le y \le 1-x^2$, $-1 \le x \le 1$. Le sens positif est : Segment $y=0$ de $A(-1, 0)$ à $B(1, 0)$, puis Parabole $y=1-x^2$ de $B(1, 0)$ à $A(-1, 0)$.

A. Calcul par intégrale curviligne
$C_1$ (Segment $y=0$ de $-1$ à $1$): $y=0$, $dy=0$.
    $$
    \int_{C_1} \vec{F}\cdot d\vec{r} = \int_{-1}^{1} y dx + (-x/2) dy = \int_{-1}^{1} 0 dx = 0
    $$
$C_2$ (Parabole $y=1-x^2$ de $1$ à $-1$): $x: 1 \to -1$. $dy=-2x dx$.
    $$
    \int_{C_2} \vec{F}\cdot d\vec{r} = \int_{1}^{-1} [y dx + (-x/2) dy] = \int_{1}^{-1} [(1-x^2) dx + (-x/2) (-2x dx)]
    $$
    $$
    = \int_{1}^{-1} [1-x^2 + x^2] dx = \int_{1}^{-1} 1 dx = [x]_{1}^{-1} = -1 - 1 = -2
    $$
$$
\int_{\mathcal{C}}\vec{F}\cdot d\vec{r} = 0 + (-2) = \mathbf{-2}
$$

B. Vérification par la Formule de Green

Rotationnel 2D :
$$
\frac{\partial Q}{\partial x} = -1/2, \quad \frac{\partial P}{\partial y} = 1
$$
$$
\frac{\partial Q}{\partial x}-\frac{\partial P}{\partial y} = -1/2 - 1 = -3/2
$$
Intégrale double :
$$
\iint_{\mathcal{R}} -3/2 dx dy = -3/2 \cdot \text{Aire}(\mathcal{R})
$$
L'aire de $\mathcal{R}$ est $\text{Aire} = \int_{-1}^{1} (1-x^2) dx = \left[x - \frac{x^3}{3}\right]_{-1}^{1} = 4/3$.
$$
\iint_{\mathcal{R}} -3/2 dx dy = -3/2 \cdot (4/3) = \mathbf{-2}
$$
Conclusion : OK
\section{Strokes}


 \end{document}
